\chapter{欠測データと EM アルゴリズム}\label{ch:16}

\begin{goalbox}
\begin{itemize}
\item 欠測メカニズム MCAR・MAR・MNAR を定義し，具体例を分類できる．
\item 完全ケース解析が偏らない条件・偏る条件を説明し，偏りの大きさを計算できる．
\item 欠測の有無と群・共変量の関連を検定し，その結果から言えること・言えないことを述べられる．
\item 観測されない度数を含む多項分布モデル（ABO 血液型）で，EM アルゴリズムの E ステップ・M ステップを導出し，反復法を構成できる．
\end{itemize}
\end{goalbox}

\begin{exambox}
\begin{itemize}
\item \kakomon{2024}{3}〔1〕：改善割合を比較する試験で，欠測（$T$群30人，$C$群10人）と群の独立性の $\chi^2$ 検定（12.5，有意）．
      同問〔2〕〜〔5〕（多重補完の統合）は，統合の式が問題文で与えられる誘導付きの設問である（\cref{app:F}）．
\item 統計数理では\kakomonSM{2016}{5}（MCAR の検定：有意でなくても MCAR とは言えない），\kakomonSM{2022}{5}（対応のあるデータで片方が欠測したときの検定）．
\item \kakomon{2019}{5}（共通）〔4〕：ABO 血液型の遺伝子頻度の推定．完全データの尤度の最大化（M ステップ）と観測されない度数の期待値（E ステップ）を誘導され，反復法を構成する．
\end{itemize}
\end{exambox}

\flowline{欠測を含む\\データ}{完全データでの\\推定対象}{完全ケース・\\尤度（EM）}{MCAR/MAR}{E・M ステップ\\の反復}{MAR は\\検証不能}

\section{問題設定}\label{sec:16-setting}
完全データを $\bm Y=(\bm Y_{\mathrm{obs}},\bm Y_{\mathrm{mis}})$，欠測の指示変数を $\bm R$（観測なら1）とする．
解析の目的は，欠測がなかった場合に推定したかった量（estimand）を，観測されたデータから偏りなく推定することである．
臨床試験では，欠測の多くは中止・脱落によって起こり，脱落は治療の効果や副作用と関連しうる．

\section{欠測メカニズム}\label{sec:16-mech}
\Hissu
\begin{teigi}[欠測メカニズム（Rubin の分類）]\label{def:16-mech}
欠測の起こり方を，欠測指示 $\bm R$ の条件付き分布によって次の3つに分類する．
\begin{itemize}
\item \textbf{MCAR}（missing completely at random）：$\Prob(\bm R\mid\bm Y_{\mathrm{obs}},\bm Y_{\mathrm{mis}})=\Prob(\bm R)$．欠測はデータと無関係．
\item \textbf{MAR}（missing at random）：$\Prob(\bm R\mid\bm Y_{\mathrm{obs}},\bm Y_{\mathrm{mis}})=\Prob(\bm R\mid\bm Y_{\mathrm{obs}})$．欠測は観測された値（共変量・過去の測定値）だけに依存．
\item \textbf{MNAR}（missing not at random）：MAR でない．すなわち，観測値を与えても欠測が観測されていない値そのものに依存する．
\end{itemize}
\end{teigi}
例：検体の破損による欠測は MCAR．前回の測定値が悪い患者が来院しなくなるのは（前回値が観測されていれば）MAR．
今回の症状が悪化したこと自体が来院しない理由なら MNAR．

\begin{meidai}[MAR の下での尤度の無視可能性]\label[meidai]{prop:16-ignor}
MAR で，データのパラメータ $\theta$ と欠測メカニズムのパラメータ $\psi$ が別個（distinct）なら，
観測データの尤度は $L(\theta,\psi)=f(\bm R\mid\bm Y_{\mathrm{obs}};\psi)\,f(\bm Y_{\mathrm{obs}};\theta)$ と分解され，
$\theta$ の推測は欠測メカニズムをモデル化せずに $f(\bm Y_{\mathrm{obs}};\theta)$ だけで行える．
\end{meidai}
\begin{proof}
$f(\bm Y_{\mathrm{obs}},\bm R)=\int f(\bm Y_{\mathrm{obs}},\bm Y_{\mathrm{mis}};\theta)f(\bm R\mid\bm Y_{\mathrm{obs}},\bm Y_{\mathrm{mis}};\psi)\,\dd\bm Y_{\mathrm{mis}}$．
MAR なら第2因子は $\bm Y_{\mathrm{mis}}$ によらず積分の外に出せて，残りは $\int f(\bm Y;\theta)\dd\bm Y_{\mathrm{mis}}=f(\bm Y_{\mathrm{obs}};\theta)$．
\end{proof}
反復測定の混合モデル（\cref{ch:08}）のような尤度に基づく解析が MAR の下で妥当なのはこのためである．
\cref{sec:16-em} の EM アルゴリズムは，この観測データの尤度 $f(\bm Y_{\mathrm{obs}};\theta)$ を最大化する計算法である．

\paragraph{データから言えること}
欠測の有無と観測された変数（群，ベースライン値）の関連を調べれば，MCAR を\textbf{否定}することはできる．
\kakomon{2024}{3}〔1〕では欠測が $T$ 群30\%，$C$ 群10\%で，$2\times2$ 表の独立性の $\chi^2$ 統計量（\cref{ch:04}）は
$\chi^2=\frac{200(30\times90-70\times10)^2}{100\cdot100\cdot40\cdot160}=12.5>\chi^2_1(0.05)=3.84$ なので，欠測は群と独立でない．
しかし関連がないことは MCAR の証明にならず（\kakomonSM{2016}{5}），MAR と MNAR は観測データからは区別できない．

\section{主な解析法}\label{sec:16-methods}
\Hissu
\begin{table}[htbp]
\centering\small
\caption{欠測への対処法と妥当となる条件}\label{tab:16-methods}
\begin{tabular}{@{}lL{48mm}L{52mm}@{}}
\toprule
方法 & 内容 & 妥当性\\
\midrule
完全ケース解析 & 欠測のある個体を除く & MCAR で不偏（精度は低下）．MAR では一般に偏る（欠測が解析モデルの共変量のみに依存する場合の回帰係数は例外）\\
単一補完（平均値，LOCF） & 欠測を1つの値で埋める & 推定対象により偏り（LOCF など），常に分散を過小評価（補完値を確定値として扱う）\\
尤度に基づく方法 & 観測データの尤度を最大化（EM など） & MAR で妥当（モデルが正しければ）\\
逆確率重み付け（IPW） & 観測確率の逆数で重み付け & MAR で妥当（観測確率のモデルが正しければ）\\
\bottomrule
\end{tabular}
\end{table}
\paragraph{完全ケース解析の偏り}
$Y$ の平均を推定したいとし，観測の有無 $R$ が共変量 $X$（全員観測）だけに依存する（MAR）とする．
完全ケースの平均の期待値は $\E[Y\mid R=1]=\sum_x\E[Y\mid X=x]\Prob(X=x\mid R=1)$ で，
$\Prob(X=x\mid R=1)\propto\Prob(R=1\mid x)\Prob(X=x)$ は $X$ の母集団分布 $\Prob(X=x)$ と異なるので，$\E[Y\mid X=x]$ が $x$ により異なれば偏る．
MCAR なら $\Prob(R=1\mid x)$ が一定なので偏らない．
\paragraph{逆確率重み付け（IPW）}
観測された個体に重み $1/\Prob(R=1\mid x)$ を付けて平均をとると，観測されやすい層の過剰代表が打ち消され，
重みの分布は $\Prob(X=x\mid R=1)/\Prob(R=1\mid x)\propto\Prob(X=x)$ となるので，MAR の下で偏りが除かれる（\cref{reidai:16-3}）．

\Youchui LOCF（最終観測値の引き継ぎ）は「脱落後の値が変わらない」という強い仮定であり，保守的とも限らない．
なお，予測分布から欠測値を複数回補完して結果を統合する\textbf{多重補完}も MAR の下で用いられるが，統合の公式は過去問では問題文で与えられており，本書では扱わない（\cref{app:F}）．
MAR は検証できない仮定なので，実務では MNAR を想定した感度分析も行う（具体的な方法は本書では扱わない）．

\section{EM アルゴリズム}\label{sec:16-em}
\Hissu
\subsection{一般形}
観測データ $\bm y$ の尤度 $L(\theta;\bm y)=f(\bm y;\theta)$ の最大化が難しいが，観測されない部分 $\bm z$ を加えた\textbf{完全データ} $(\bm y,\bm z)$ の対数尤度 $\ell_c(\theta;\bm y,\bm z)$ なら容易に最大化できるとき，
現在の値 $\theta^{(t)}$ から次の2ステップを交互に繰り返す．
\begin{itemize}
\item \textbf{E ステップ}：観測データ $\bm y$ を条件とした期待値 $Q(\theta\mid\theta^{(t)})=\E_{\theta^{(t)}}\bigl[\ell_c(\theta;\bm y,\bm Z)\bigm|\bm y\bigr]$ を計算する．
\item \textbf{M ステップ}：$\theta^{(t+1)}=\arg\max_\theta Q(\theta\mid\theta^{(t)})$．
\end{itemize}
$\ell_c$ が観測されない度数について線形（多項分布など）なら，E ステップは「観測されない度数を，観測データを条件とした期待値で置き換える」ことと同じである．
この標準的な EM では，各反復で観測データの尤度は減少しない：$L(\theta^{(t+1)};\bm y)\ge L(\theta^{(t)};\bm y)$（証明は省略）．
この性質は，E ステップの期待値が\textbf{観測データ $\bm y$ を条件とした}ものであることに依存する．

\subsection{例：ABO 血液型の遺伝子頻度}
遺伝子 O，A，B の頻度を $r,p,q$（$r+p+q=1$）とし，ハーディ・ワインベルグ平衡で遺伝子型 OO, AA, AO, BB, BO, AB の確率が
$r^2,p^2,2pr,q^2,2qr,2pq$ とする．観測されるのは表現型の度数 $n_O,n_A,n_B,n_{AB}$（合計 $N$）だけで，
$f_{OO}=n_O$，$f_{AA}+f_{AO}=n_A$，$f_{BB}+f_{BO}=n_B$，$f_{AB}=n_{AB}$ のうち，A型・B型の内訳 $f_{AA},f_{AO},f_{BB},f_{BO}$ が観測されない．

\paragraph{M ステップ（完全データの最尤推定）}
完全データの対数尤度は，定数を除いて
\[
\ell_c=(2f_{OO}+f_{AO}+f_{BO})\log r+(2f_{AA}+f_{AO}+f_{AB})\log p+(2f_{BB}+f_{BO}+f_{AB})\log q .
\]
制約 $r+p+q=1$ の下で，ラグランジュ乗数 $\lambda$ を用いて $\ell_c-\lambda(r+p+q-1)$ を偏微分して0とおくと，
$2f_{OO}+f_{AO}+f_{BO}=\lambda r$ などの3式を得る．3式を足すと左辺は全遺伝子数 $2N$ なので $\lambda=2N$．よって
\begin{equation}
\hat r=\frac{2f_{OO}+f_{AO}+f_{BO}}{2N},\qquad
\hat p=\frac{2f_{AA}+f_{AO}+f_{AB}}{2N},\qquad
\hat q=\frac{2f_{BB}+f_{BO}+f_{AB}}{2N}\label{eq:16-mstep}
\end{equation}
（遺伝子の数え上げ）．
\paragraph{E ステップ（観測データを条件とした期待値）}
A型の各人は，独立に確率 $\dfrac{p^2}{p^2+2pr}$ で AA，残りの確率で AO である（条件付き確率 $\Prob(\text{AA}\mid\text{A型})$）．
よって $n_A$ を与えると $f_{AA}\mid n_A\sim\Bin\bigl(n_A,\ p^2/(p^2+2pr)\bigr)$ で，
\begin{equation}
\E[f_{AA}\mid n_A]=n_A\frac{p^2}{p^2+2pr}=n_A\frac{p}{p+2r},\qquad \E[f_{AO}\mid n_A]=n_A\frac{2r}{p+2r},\label{eq:16-estep}
\end{equation}
B型も同様に $\E[f_{BB}\mid n_B]=n_B\,q/(q+2r)$，$\E[f_{BO}\mid n_B]=n_B\,2r/(q+2r)$．
すなわち，観測された A型の人数を，現在のパラメータでの AA と AO の条件付き確率で按分する．
\paragraph{反復法}
初期値 $(r^{(0)},p^{(0)},q^{(0)})$ から，\cref{eq:16-estep} に $(r^{(t)},p^{(t)},q^{(t)})$ を代入して $f_{AA}^{(t)},f_{AO}^{(t)},f_{BB}^{(t)},f_{BO}^{(t)}$ を求め（E ステップ），
それを\cref{eq:16-mstep} に代入して $(r^{(t+1)},p^{(t+1)},q^{(t+1)})$ を求める（M ステップ），を変化が十分小さくなるまで繰り返す．
按分なので $f_{AA}^{(t)}+f_{AO}^{(t)}=n_A$ が常に成り立ち，$r+p+q=1$ も保たれる．

\begin{warnbox}[試験の公式解答との相違]
\kakomon{2019}{5}〔4〕(ii) の公式解答（解答例）は，(i) で求めた $p,q$ を用いた観測されない度数の期待値として
$\E[f_{AA}]=Np^2$，$\E[f_{AO}]=n_A-Np^2$（B型も同様に $Nq^2$，$n_B-Nq^2$）を示し，これと (i) の式（\cref{eq:16-mstep}）を交互に用いる反復を EM アルゴリズムとして示している．
また，$f_{AO}=N\cdot2pr$ とする反復は，$f_{AA}+f_{AO}=n_A$ が成り立つとは限らないのでうまくいかない，と注記している．
\begin{enumerate}[label=(\arabic*)]
\item \textbf{両者は一般に同じ反復にならない}．$Np^2$ は観測値 $n_A$ を条件としない期待値であり，
      標準的な E ステップ $n_Ap^2/(p^2+2pr)$（\cref{eq:16-estep}）とは $n_A=N(p^2+2pr)$ のときを除いて異なる．
      収束先も一般には異なる：公式解答の反復の不動点は $2Np=Np^2+n_A+n_{AB}$，すなわち $(1-p)^2=(n_O+n_B)/N$ を満たし，
      これは観測データの最尤推定値とは一般に一致しない（\cref{reidai:16-4}(3)）．
\item \textbf{「各反復で観測データの尤度が減少しない」は標準 EM についてのみ成り立つ}．公式解答の反復にはこの保証はなく，MLE への収束も保証されない．
\item \textbf{試験での答え方}：問題文の誘導（(i) を M ステップ，(ii) を E ステップとして交互に反復する枠組み）には必ず従う．
      (ii) の期待値は，何を条件とした期待値かを明記して書く．観測値 $n_A,n_B$ を条件とした $n_Ap^2/(p^2+2pr)$ などが標準的な EM の E ステップであり，
      採点基準は公表されていないので，標準形で書いたうえで，公式解答の形 $Np^2$（観測値を条件としない期待値）とは異なることを一言添えておくとよい．
\end{enumerate}
\end{warnbox}

\section{仮定}\label{sec:16-assump}
\begin{itemize}
\item 完全ケース解析：MCAR（またはそれに準じる条件）．
\item 尤度法（EM を含む）・IPW：MAR（検証不能），モデルの正しさ（IPW では観測確率のモデル）．
\item ABO の EM：ハーディ・ワインベルグ平衡（遺伝子型の確率が遺伝子頻度の積），個体の独立性．
\end{itemize}

\section{結果の解釈}\label{sec:16-interp}
\begin{itemize}
\item 群によって欠測割合が異なる（\kakomon{2024}{3}）なら，完全ケース解析は偏る可能性が高い．欠測の理由（有害事象，効果不十分）を群ごとに確認する．
\item 欠測と観測変数の関連が有意でなくても MCAR の証明にはならない．MAR か MNAR かはデータから判定できない．
\item EM の収束値は観測データの尤度の（局所）最大点である．複数の初期値から始めて同じ値に収束するかを確かめるとよい．
\end{itemize}

\section{手法選択}\label{sec:16-choice}
\begin{itemize}
\item 欠測が MCAR とみなせる → 完全ケース解析（不偏だが精度は低下）．
\item 欠測が観測された共変量・過去の値に依存（MAR）→ 尤度に基づく方法（混合モデル（\cref{ch:08}），EM）または IPW．
\item 観測されない度数・潜在的なクラスを含む多項モデルで MLE を求める → EM アルゴリズム．
\end{itemize}

\section{よくある誤り}\label{sec:16-err}
\begin{warnbox}
\begin{itemize}
\item 欠測と群・共変量の関連が有意でないことを MCAR の証明とする．
\item 単一補完したデータを完全データとして分析し，標準誤差を過小評価する．
\item MAR を「欠測がランダム（データと無関係）」と読み違える（それは MCAR）．
\item EM の E ステップで，観測値を条件としない期待値（$Np^2$ など）を用いて，標準 EM の性質（尤度の単調性・MLE への収束）が成り立つとする．
\end{itemize}
\end{warnbox}

\section{数理統計との接続}\label{sec:16-link}
\begin{linkbox}
\begin{itemize}
\item E ステップの $f_{AA}\mid n_A\sim\Bin(n_A,p^2/(p^2+2pr))$ は，多項分布の一部のセルを併合したときの条件付き分布（多項分布の条件付き分布は多項分布）である．
\item M ステップは多項分布の MLE（セルの相対度数）を遺伝子単位で数えたものである．
\item \kakomonSM{2018}{3}（0切断二項分布の反復解法）は，観測されないゼロの個数を補う EM アルゴリズムの更新と一致する．
\end{itemize}
\end{linkbox}

\section{例題}\label{sec:16-ex}
\begin{reidai}[欠測メカニズムの分類]\label{reidai:16-1}
次の欠測を MCAR・MAR・MNAR に分類せよ（必要な仮定があれば述べよ）．
\begin{enumerate}[label=(\arabic*)]
\item 検査機器の故障で，ある日の全患者の測定値が欠測した．
\item 第4週の血圧が高かった患者ほど第8週に来院しなかった（第4週の値は観測済み）．
\item 抑うつ尺度の評価で，症状が最も悪い時期の患者が回答を拒否した．
\end{enumerate}
\end{reidai}

\begin{reidai}[完全ケース解析の偏り]\label{reidai:16-3}
2 値の共変量 $X$（各50\%）で $\E[Y\mid X=0]=10$，$\E[Y\mid X=1]=20$．$Y$ の観測確率は $X=0$ で0.9，$X=1$ で0.5（$X$ は全員観測）．
\begin{enumerate}[label=(\arabic*)]
\item 欠測メカニズムは何か．
\item 完全ケースでの $Y$ の平均の期待値を求め，真の平均15と比べよ．
\item 観測確率の逆数で重み付けした平均の期待値を求めよ．
\end{enumerate}
\end{reidai}

\begin{reidai}[ABO 血液型の EM アルゴリズム]\label{reidai:16-4}
$N=200$ 人の表現型が $n_O=60$，$n_A=80$，$n_B=40$，$n_{AB}=20$ であった．初期値を $(r^{(0)},p^{(0)},q^{(0)})=(0.5,\,0.25,\,0.25)$ とする．
\begin{enumerate}[label=(\arabic*)]
\item 標準的な EM の E ステップで $f_{AA},f_{AO},f_{BB},f_{BO}$ を求め，M ステップで $(r^{(1)},p^{(1)},q^{(1)})$ を求めよ．
\item E ステップの代わりに $f_{AA}=Np^2$，$f_{AO}=n_A-Np^2$（B型も同様）を用いた場合の $(r^{(1)},p^{(1)},q^{(1)})$ を求めよ．
\item (2) の反復が収束するとき，その極限の $p,q$ を求めよ（$\sqrt{0.5}=0.7071$，$\sqrt{0.7}=0.8367$）．
\end{enumerate}
\end{reidai}

\section{解答}\label{sec:16-sol}
\begin{kaitou}{reidai:16-1}
(1) MCAR（故障が患者の状態と無関係なら）．
(2) MAR（欠測が観測済みの第4週の値に依存し，それを与えれば第8週の値によらないなら）．第8週の値自体に依存するなら MNAR．
(3) MNAR（欠測が欠測値そのもの＝その時点の症状に依存）．
\end{kaitou}

\begin{kaitou}{reidai:16-3}
(1) 欠測は観測された $X$ だけに依存するので MAR．\\
(2) 観測者の構成は $X=0$：$0.5\times0.9=0.45$，$X=1$：$0.5\times0.5=0.25$．平均 $\frac{0.45\times10+0.25\times20}{0.70}=\frac{9.5}{0.7}=13.57$．真の15より小さく偏る．\\
(3) 重み $1/0.9$，$1/0.5$ で $\frac{0.45\times10/0.9+0.25\times20/0.5}{0.45/0.9+0.25/0.5}=\frac{5+10}{1}=15$．偏りが除かれる．
（$X$ ごとの平均を $X$ の分布で平均すること，すなわち $X$ を条件とした解析でも同じ結果になる．）
\end{kaitou}

\begin{kaitou}{reidai:16-4}
(1) $p/(p+2r)=0.25/1.25=0.2$，$q/(q+2r)=0.2$ なので
$f_{AA}=80\times0.2=16$，$f_{AO}=64$，$f_{BB}=40\times0.2=8$，$f_{BO}=32$．\cref{eq:16-mstep}より
\[
r^{(1)}=\frac{120+64+32}{400}=0.54,\qquad p^{(1)}=\frac{32+64+20}{400}=0.29,\qquad q^{(1)}=\frac{16+32+20}{400}=0.17 .
\]
(2) $f_{AA}=200\times0.0625=12.5$，$f_{AO}=67.5$，$f_{BB}=12.5$，$f_{BO}=27.5$ で
\begin{gather*}
r^{(1)}=\frac{120+67.5+27.5}{400}=0.5375,\qquad p^{(1)}=\frac{25+67.5+20}{400}=0.28125,\\
q^{(1)}=\frac{25+27.5+20}{400}=0.18125 .
\end{gather*}
同じ初期値から異なる値になり，両者は別の反復である．\\
(3) 不動点では $p=\{2Np^2+(n_A-Np^2)+n_{AB}\}/(2N)$，すなわち $Np^2-2Np+(n_A+n_{AB})=0$，$(1-p)^2=(n_O+n_B)/N=100/200=0.5$．
$0<p<1$ より $p=1-\sqrt{0.5}=0.293$．同様に $(1-q)^2=(n_O+n_A)/N=0.7$ から $q=1-\sqrt{0.7}=0.163$．
これは観測データの最尤推定値（標準 EM の収束値）とは一般に一致しない
（参考：この数値で標準 EM を計算機で反復すると $\hat p=0.2923$ に収束し，$1-\sqrt{0.5}=0.2929$ とわずかに異なる）．
\end{kaitou}

\begin{summarybox}
\begin{itemize}
\item MCAR：データと無関係．MAR：観測値のみに依存．MNAR：欠測値自体に依存．MAR と MNAR はデータで区別できない．
\item 欠測と観測変数の関連の検定で MCAR は否定できるが，非有意でも MCAR の証明にはならない．
\item MAR＋パラメータの分離で，尤度は欠測メカニズムを無視できる．完全ケースは MCAR で不偏，MAR では一般に偏る（IPW で補正できる）．単一補完は分散を過小評価．
\item EM：E ステップ＝観測データを条件とした完全データ対数尤度の期待値，M ステップ＝その最大化．標準 EM では観測データの尤度は減少しない．
\item ABO：E ステップ $\E[f_{AA}\mid n_A]=n_Ap^2/(p^2+2pr)$，M ステップ $\hat p=(2f_{AA}+f_{AO}+f_{AB})/(2N)$ など（遺伝子の数え上げ）．
\item \kakomon{2019}{5} の公式解答の $\E[f_{AA}]=Np^2$ は観測値を条件としない期待値で，標準 EM とは別の反復になる．
\end{itemize}
\end{summarybox}
